\documentclass[10pt,letterpaper]{article}
% Shared corporate style for StrategicDefense Systems Security LTD.
\usepackage[T1]{fontenc}
\usepackage[utf8]{inputenc}
\usepackage{tgtermes}
\usepackage{microtype}
\usepackage{xcolor}
\usepackage{graphicx}
\usepackage{geometry}
\usepackage{fancyhdr}
\usepackage{lastpage}
\usepackage{enumitem}
\usepackage{tabularx}
\usepackage{array}
\usepackage{booktabs}
\usepackage{setspace}
\usepackage{ragged2e}
\usepackage[hidelinks]{hyperref}

\definecolor{SDSSNavy}{HTML}{173C64}
\definecolor{SDSSGreen}{HTML}{183C31}
\definecolor{SDSSRule}{HTML}{A9B8C8}

\newcommand{\CompanyFull}{StrategicDefense Systems Security LTD.}
\newcommand{\CompanyShort}{Strategic Defense}
\newcommand{\ContactLine}{4 Oguda Street, AMAC, Abuja, FCT, Nigeria\enspace\textbar\enspace Email: info@strategic-defence.com\enspace\textbar\enspace Telephone: +234 703 867 6619}
\graphicspath{{../../}}

\hypersetup{
  colorlinks=true,
  linkcolor=SDSSNavy,
  urlcolor=SDSSNavy,
  pdfauthor={StrategicDefense Systems Security LTD.}
}

\newcommand{\SDSSLogo}[1][2.6in]{\includegraphics[width=#1]{logo_transparent.png}}
\newcommand{\SDSSIcon}[1][0.15in]{\includegraphics[width=#1]{icon 1.png}}

\geometry{left=0.72in,right=0.72in,top=0.88in,bottom=0.68in,headheight=47pt,headsep=12pt,footskip=25pt}
\setlength{\parindent}{0pt}
\setlength{\parskip}{5pt}
\setstretch{1.02}
\setcounter{tocdepth}{2}

\hypersetup{
  pdftitle={PRELIMINARY TECHNICAL PROPOSAL - SHIELD},
  pdfsubject={Secure Homeland Integrated Early-Response and Land-Border Defence}
}

\makeatletter
\renewcommand\section{\@startsection{section}{1}{0pt}%
  {-2.3ex plus -0.8ex minus -0.2ex}{0.9ex plus 0.2ex}%
  {\color{SDSSNavy}\large\bfseries}}
\renewcommand\subsection{\@startsection{subsection}{2}{0pt}%
  {-1.8ex plus -0.6ex minus -0.2ex}{0.6ex plus 0.2ex}%
  {\color{SDSSNavy}\normalsize\bfseries}}
\makeatother
\renewcommand{\contentsname}{Contents}

\fancypagestyle{sdss}{%
  \fancyhf{}%
  \fancyhead[L]{\raisebox{-4pt}{\SDSSLogo[1.45in]}}%
  \fancyfoot[L]{\raisebox{-1pt}{\SDSSIcon}\hspace{5pt}{\color{SDSSNavy}\footnotesize\CompanyFull}}%
  \fancyfoot[C]{\footnotesize\thepage\ / \pageref{LastPage}}%
  \renewcommand{\headrulewidth}{0.7pt}%
  \renewcommand{\footrulewidth}{0.5pt}%
  \renewcommand{\headrule}{\hbox to\headwidth{\color{SDSSNavy}\leaders\hrule height \headrulewidth\hfill}}%
  \renewcommand{\footrule}{\hbox to\headwidth{\color{SDSSNavy}\leaders\hrule height \footrulewidth\hfill}}%
}

\newcommand{\EquipmentItem}[2]{\item[\textbf{#1}] #2}
\newcommand{\PrincipleItem}[2]{\item \textbf{#1:} #2}

\begin{document}
\hypersetup{pageanchor=false}

% Cover page
\thispagestyle{empty}
\SDSSLogo[2.65in]

\vspace{0.52in}
{\color{SDSSNavy}\rule{\textwidth}{0.8pt}}

\vspace{0.18in}
{\color{SDSSNavy}\fontsize{23}{27}\selectfont\bfseries PRELIMINARY TECHNICAL PROPOSAL\par}
\vspace{0.16in}
{\color{SDSSNavy}\fontsize{19}{23}\selectfont\bfseries SHIELD\par}
\vspace{0.06in}
{\color{SDSSNavy}\fontsize{12.5}{16}\selectfont\bfseries SECURE HOMELAND INTEGRATED EARLY-RESPONSE AND LAND-BORDER DEFENCE\par}

\vspace{0.18in}
{\color{SDSSNavy}\rule{\textwidth}{0.8pt}}

\vfill
\renewcommand{\arraystretch}{1.65}
\begin{tabular}{@{}l@{}l@{}}
  \makebox[1.85in][l]{\color{SDSSNavy}\bfseries Prepared by} & \CompanyFull{} (\CompanyShort) \\
  \makebox[1.85in][l]{\color{SDSSNavy}\bfseries Document status} & Preliminary for Technical Discussion \\
  \makebox[1.85in][l]{\color{SDSSNavy}\bfseries Version} & 1.0 \\
  \makebox[1.85in][l]{\color{SDSSNavy}\bfseries Date} & 26 September 2026 \\
  \makebox[1.85in][l]{\color{SDSSNavy}\bfseries Document reference} & SDSS/SHIELD/PTP/001 \\
\end{tabular}

\vfill
\begin{tabular*}{\textwidth}{@{\extracolsep{\fill}}lr}
  {\footnotesize\CompanyFull} & {\footnotesize Commercial in Confidence}
\end{tabular*}
\clearpage

% Contents
\hypersetup{pageanchor=true}
\pagenumbering{roman}
\pagestyle{sdss}
{\color{SDSSNavy}\tableofcontents}
\clearpage

% Main matter
\pagenumbering{arabic}
\pagestyle{sdss}
\thispagestyle{sdss}

\section{PURPOSE OF THIS PROPOSAL}

This preliminary proposal presents the SHIELD concept and the equipment \CompanyShort{} offers to supply. It is a basis for technical discussion with the Office of the Special Adviser, not a final offer.

The final proposal will follow a site survey. It will confirm exact quantities, locations, and schedule.

\section{BACKGROUND AND OUR UNDERSTANDING}

Nigeria's security agencies already hold a wealth of information and experience. The opportunity now is to move that information more quickly from collection to action, so that agencies in the field can act on it together.

The Office of the Special Adviser to the President on Homeland Security, created in May 2026, exists to close that gap through intelligence integration, inter-agency coordination, and proactive risk management. It builds on the Intelligence Fusion Centre established under the Office of the National Security Adviser, which already brings the major security agencies together in Abuja.

The next step is to extend that coordination to where threats occur: the land borders and the scene of a crisis. Two practical gaps stand out:

\begin{enumerate}[label=\alph*.,leftmargin=2em,itemsep=5pt]
  \item \textbf{Unwatched borders.} Long stretches of land border have no persistent surveillance, so illegal crossings go undetected.
  \item \textbf{Separate radio systems.} Police, military and paramilitary units often operate on separate radio systems, which can make direct communication at an incident difficult.
\end{enumerate}

Neither gap can be closed by the military alone. Nigeria's 2019 National Security Strategy defines national security broadly, reaching beyond military matters to economic, social, environmental and human-rights concerns. Its list of national threats includes porous borders, transnational organised crime, banditry and kidnapping, and public health challenges. Meeting them is shared work: the police, customs, immigration, civil defence, emergency services and communities each play a part alongside the armed forces. What SHIELD adds is earlier, clearer information where it is needed, at the border and at the scene: knowing who is moving, where, when and how, early enough to act before harm is done.

SHIELD addresses both with equipment, delivered and working, not with further studies.

\section{THE SHIELD CONCEPT}

SHIELD is built on Nigeria's National Security Strategy 2019 (NSS), the overarching framework for national security. The NSS takes a comprehensive view of security centred on the well-being of people, recognises that security is no longer the exclusive preserve of the military, and calls for a coordinated, whole-of-nation approach. SHIELD equips civilian and military agencies to deliver on five of its commitments.

\renewcommand{\arraystretch}{1.22}
\begin{tabularx}{\textwidth}{@{}>{\RaggedRight\arraybackslash}p{1.55in}>{\RaggedRight\arraybackslash}p{3.10in}>{\RaggedRight\arraybackslash}X@{}}
\toprule
\textcolor{SDSSNavy}{\textbf{Commitment}} & \textcolor{SDSSNavy}{\textbf{NSS}} & \textcolor{SDSSNavy}{\textbf{SHIELD}} \\
\midrule
\textcolor{SDSSNavy}{\textbf{Porous borders}} & Nigeria's land borders are largely porous and inadequately manned; the South-West border with Benin is known for smuggling; stronger border security is critical to countering terrorism and organised crime & Package A - continuous monitoring of a pilot section of the South-West border \\
\textcolor{SDSSNavy}{\textbf{Transnational organised crime}} & Smuggling, the spread of small arms and trafficking are fuelled in part by weak border security management & Package A - sensors, cameras and cargo scanning at border posts \\
\textcolor{SDSSNavy}{\textbf{Internal security}} & Commits to stronger border security and to technology-enabled intelligence infrastructure & Package A - a joint border control room shared by all border agencies \\
\textcolor{SDSSNavy}{\textbf{Crisis response and resilience}} & Calls for early warning and response, effective communication, and inter-agency crisis teams with unified command, control, communication and coordination & Package B - radio bridging, a mobile incident command vehicle and emergency communications \\
\textcolor{SDSSNavy}{\textbf{Environmental threats}} & Treats floods and other natural and man-made disasters as national security threats needing contingency plans & Package B - coordination and communications that keep working when normal networks fail \\
\bottomrule
\end{tabularx}

SHIELD also supports the Secure and Prepare work streams of the National Counter-Terrorism Strategy (NACTEST), on which the NSS relies for its counter-terrorism aims.

Two principles guide the design:

\begin{enumerate}[label=\alph*.,leftmargin=2em,itemsep=5pt]
  \item \textbf{Shared picture, not shared command.} Each agency keeps its mandate and chain of command. SHIELD gives them the means to see and talk to each other.
  \item \textbf{Build on what exists.} SHIELD extends the existing coordination structure rather than creating a new one.
\end{enumerate}

\section{PACKAGE A - BORDER MONITORING AND MANAGEMENT}

Package A gives the agencies that manage Nigeria's land borders a clear, continuous view of what crosses them. The main users are the Nigeria Customs Service, the Nigeria Immigration Service, the Police and the Nigeria Security and Civil Defence Corps, working alongside the military where needed.

The aim is not only to intercept armed threats. It is also to keep border communities safe in their daily lives, curb smuggling and human trafficking, and let legitimate trade and travellers move more safely and quickly.

The stakes are high. The National Security Strategy, citing the Small Arms Survey, estimates that some one million small arms and light weapons are in circulation in Nigeria, and links their spread to conflicts in neighbouring countries and weak border security management.

The pilot covers a section of the South-West border with Benin, one of Nigeria's busiest trade corridors, and expands to other sections once proven. The exact section, for example along the main road corridors and the bush crossings between them, will be fixed during the site survey.

\subsection{How it works}

\begin{enumerate}[leftmargin=2em,itemsep=4pt]
  \item \textbf{Detect:} radar and ground sensors pick up people and vehicles crossing away from official posts, day and night.
  \item \textbf{Verify:} the nearest camera turns toward the alert, so an operator can see what it is before anyone is sent, which avoids wasted patrols and false alarms.
  \item \textbf{Share:} video and alerts reach the control room, where all border agencies work from the same picture, and on to national level.
  \item \textbf{Respond:} the agency best placed to deal with the situation acts, whether a customs team, an immigration officer, a police unit, or a military patrol for armed threats.
\end{enumerate}

\subsection{Equipment}

\begin{description}[style=nextline,leftmargin=0pt,labelsep=0pt,itemsep=6pt]
  \EquipmentItem{Monitoring towers}{Masts with day, night and thermal cameras and detection radar, covering open and bush terrain.}
  \EquipmentItem{Ground sensors}{Buried sensors along known unofficial crossing routes that alert when someone passes.}
  \EquipmentItem{Border post upgrades}{CCTV, vehicle barriers and gates, and space, power and network ready for ID devices.}
  \EquipmentItem{Joint border control room}{A room near the border where staff from all border agencies monitor cameras and sensors together.}
  \EquipmentItem{Tower power}{Solar panels with battery storage, independent of the grid.}
  \EquipmentItem{Communications links}{Radio and satellite links carrying video and alerts to the control room and onward.}
  \EquipmentItem{Mobile monitoring units}{Vehicles with a camera and radar mast, moved to where crossing patterns shift.}
  \EquipmentItem{Portable observation kits}{Portable day and night cameras carried by border officers to cover gaps between towers.}
  \EquipmentItem{Vehicle and cargo scanning}{X-ray scanners so customs can inspect vehicles and cargo for contraband without unloading, speeding up legitimate trade.}
  \EquipmentItem{Site protection}{Fencing and controlled entry around border posts and the control room.}
\end{description}

\section{PACKAGE B - JOINT EMERGENCY RESPONSE AND COMMUNICATIONS}

Package B helps every agency that responds to an emergency work as one team at the scene. That includes the police, civil defence, fire service, the National Emergency Management Agency (NEMA) and state emergency agencies, the Federal Road Safety Corps, health services and, where needed, the military.

Emergencies include not only attacks and kidnappings but also floods, building collapses, chemical spills and other hazardous-material incidents, major road accidents, fires and large public events.

In June 2026, NEMA and the World Bank began reviewing Nigeria's disaster management frameworks, with the aim of strengthening coordination between agencies, clarifying roles and closing gaps in response structures. Package B supports that effort with the equipment that lets agencies work as one team at the scene.

\subsection{How it works}

At any major incident, responders arrive on their own radio networks and cannot talk to each other. The radio bridging equipment links those networks, so the incident commander can open a shared channel for all responders in minutes, without anyone changing radios.

The mobile incident command vehicle becomes the on-scene coordination point. Agency leads plan together inside it, with a satellite link to national level. If phone networks are down, as often happens in floods or large incidents, the emergency communication kits keep the scene connected.

\subsection{Equipment}

\begin{description}[style=nextline,leftmargin=0pt,labelsep=0pt,itemsep=6pt]
  \EquipmentItem{Radio bridging equipment}{Gateways linking the radio systems of different agencies, fixed at coordination centres and portable for the field.}
  \EquipmentItem{Mobile incident command vehicle}{One equipped vehicle with radios, bridging equipment, satellite link, displays and a meeting space for agency leads.}
  \EquipmentItem{Emergency communication kits}{Satellite phones and portable satellite terminals for when phone and internet networks fail.}
  \EquipmentItem{Deployable coordination shelters}{Quick-erect shelters with lighting and power for longer operations or places a vehicle cannot reach.}
  \EquipmentItem{Standby power}{Generators and battery packs for the vehicle, shelters and kits.}
  \EquipmentItem{Mobile relay masts}{Portable masts with secure radio relay, extending coverage into remote areas.}
\end{description}

\clearpage
\section{DESIGN PRINCIPLES}
\thispagestyle{sdss}

All SHIELD equipment is chosen against six principles:

\begin{enumerate}[label=\alph*.,leftmargin=2em,itemsep=5pt]
  \PrincipleItem{Technology-neutral}{no brand is fixed in advance; equipment is selected against performance requirements.}
  \PrincipleItem{Open standards}{video, data, and radio interfaces use recognised standards, so SHIELD connects to national systems and avoids lock-in to a single supplier.}
  \PrincipleItem{Modular and scalable}{the pilot section and first vehicle form a template that can be copied to other border sections and zones.}
  \PrincipleItem{Built for the environment}{equipment is rated for heat, humidity, heavy rain, and dust.}
  \PrincipleItem{Independent power}{towers and field sites run on solar and batteries; command equipment carries its own backup.}
  \PrincipleItem{Protected against tampering and cyber attack}{links are encrypted, equipment is physically hardened, and access is controlled.}
\end{enumerate}

\section{IMPLEMENTATION APPROACH}

SHIELD is delivered in steps, starting with an early live demonstration of radio bridging with participating agencies. Timings for each step will be confirmed in the final proposal, after the site survey.

\begin{description}[style=nextline,leftmargin=0pt,labelsep=0pt,itemsep=7pt]
  \EquipmentItem{Step 0 - Site survey and design}{Survey the pilot border section and coordination sites; confirm the final equipment list and installation design.}
  \EquipmentItem{Step 1 - Launch and demonstration}{Equipment ordered and export licenses applied for; live demonstration of radio bridging with participating agencies.}
  \EquipmentItem{Step 2 - Interoperable communications}{Radio bridging equipment and emergency communication kits delivered, installed, and handed over.}
  \EquipmentItem{Step 3 - Incident command vehicle and border pilot}{Mobile incident command vehicle delivered; pilot border section installed and operational.}
  \EquipmentItem{Step 4 - Expansion}{Further border sections and response units, based on pilot results.}
\end{description}

\section{TRAINING AND SUSTAINMENT}

\CompanyShort{} trains the people who use SHIELD and keeps the equipment running.

\begin{enumerate}[label=\alph*.,leftmargin=2em,itemsep=5pt]
  \item \textbf{Operator training:} control room staff, vehicle crews, and radio operators, on the equipment they use.
  \item \textbf{Technician training:} first-line maintenance and fault-finding for local technical staff.
  \item \textbf{Train-the-trainer:} Nigerian instructors certified to train new personnel.
  \item \textbf{Maintenance and spares:} a multi-year support period with preventive maintenance, a spare parts pool, and local technical staff.
\end{enumerate}

\section{WHY \MakeUppercase{\CompanyShort}}

\CompanyShort{} is a Nigerian-registered company based in Abuja, with an NCAGE code, supplying and integrating security systems for police, customs, immigration, emergency services, paramilitary services, armed forces, and government agencies.

\begin{enumerate}[label=\alph*.,leftmargin=2em,itemsep=6pt]
  \item \textbf{Relevant capability:} border and perimeter surveillance, vehicle-mounted and portable observation systems, secure communications and relay masts, command vehicle fit-out, checkpoint screening, and command-centre equipment.
  \item \textbf{Established manufacturer network:} \CompanyShort{} represents leading international manufacturers under executed agreements, and selects the system that suits the operating environment rather than a single brand.
  \item \textbf{One point of accountability:} site survey, procurement and export documentation, shipping and customs, installation and integration with existing cameras and radios, training, and through-life support, all under one contract.
  \item \textbf{Knowledge of Nigerian institutions:} engagements with the Nigerian Air Force mean SHIELD is designed around how Nigerian agencies actually operate, and can link with air surveillance assets where the client wishes.
\end{enumerate}

\section{NEXT STEPS}

\CompanyShort{} proposes three steps toward a final proposal:

\begin{enumerate}[leftmargin=2em,itemsep=6pt]
  \item A meeting with the Office of the Special Adviser to present SHIELD and answer questions.
  \item Agreement on the pilot border section and the participating agencies.
  \item A site survey, leading to the final technical proposal with confirmed quantities, locations, and schedule.
\end{enumerate}

\vfill
\begin{center}
  {\color{SDSSNavy}\rule{0.7\textwidth}{0.6pt}}\\[7pt]
  {\color{SDSSNavy}\bfseries End of Preliminary Technical Proposal SDSS/SHIELD/PTP/001 v1.0}
\end{center}

\end{document}
